\problemname{Wildfire Evacuation}

A wildfire has broken out across a region with several towns. Each town has a
number of people who need to evacuate, and each has to travel along a network
of roads to reach one of several emergency shelters. Every road has a maximum 
number of people who can travel along it, and every shelter has a maximum number 
of people it can safely house. As the coordinator overseeing the evacuation, your 
job is to figure out: what is the maximum number of people who can be moved to 
safety, given these constraints?

\section*{Input}
Towns and shelters share one numbering scheme: towns are numbered
$1$ to $T$, and shelters are numbered $T{+}1$ to $T{+}S$.

\begin{itemize}
\item The first line contains three integers $T$, $S$, $R$
  ($1 \le T, S \le 200$, $1 \le R \le 2000$): the number of towns,
  shelters, and roads.
\item The second line contains $T$ integers $p_1, \dots, p_T$
  ($0 \le p_i \le 10^9$): the number of people needing evacuation
  from each town.
\item The third line contains $S$ integers $c_1, \dots, c_S$
  ($0 \le c_i \le 10^9$): the capacity of each shelter.
\item Each of the next $R$ lines contains three integers $u$, $v$, $w$
  ($1 \le u, v \le T+S$, $1 \le w \le 10^9$): a one-way road from
  location $u$ to location $v$ with capacity $w$.
\end{itemize}

People may pass through any town or shelter on their way to the shelter
where they stay; only people who remain at a shelter count towards its
capacity. There may be more than one road between the same pair of
locations, and a road may start and end at the same location.

\section*{Output}
Output a single integer: the maximum total number of people that can be
evacuated to shelters.